\subsection{Validación de la adquisición y del montaje experimental}
\label{subsec:validacion_adquisicion_montaje}

La primera parte de la validación se centró en comprobar que la cadena baja de adquisición funcionaba de forma continua antes de interpretar la señal desde el punto de vista espectral o musical. La ruta validada fue:

\begin{center}
\texttt{ADS1299 $\rightarrow$ SPI/RDATAC $\rightarrow$ firmware $\rightarrow$ Bridge $\rightarrow$ Python $\rightarrow$ CSV}
\end{center}

El firmware reconstruye las tramas RDATAC de 24 bits, valida el prefijo de estado del ADS1299, convierte las cuentas del ADC a microvoltios y envía al backend bloques \texttt{eeg\_block\_uV} de 8 muestras. Con una frecuencia de muestreo de \SI{250}{\hertz}, cada bloque representa \SI{32}{\milli\second} de señal. Esta comprobación permite separar los problemas propios de adquisición y transporte de los problemas asociados al montaje bioeléctrico, como ruido de red, movimiento de electrodos, impedancia de contacto o artefactos fisiológicos.

Como prueba diagnóstica de la ruta ADC y del transporte digital se utilizó una captura con entradas cortocircuitadas, identificada como \texttt{20260523-175959\_post\_configfix\_shorted\_inputs}. Esta prueba pretende aislar el comportamiento del ADS1299 y del flujo de datos, sin representar una señal EEG real. El resultado fue una frecuencia efectiva de \SI{250}{\hertz}, sin discontinuidades temporales y sin estados ADS inválidos. Además, la señal presentó un valor RMS de \SI{0.115}{\micro\volt}, coherente con una entrada cortocircuitada y con bajo ruido interno.

\begin{table}[!htbp]
    \centering
    \caption{Resultado de la prueba diagnóstica de adquisición con entradas cortocircuitadas.}
    \label{tab:validacion_shorted_inputs}
    \begin{tabular}{l c}
        \hline
        \textbf{Métrica} & \textbf{Valor} \\
        \hline
        Frecuencia efectiva & \SI{250.0}{\hertz} \\
        Discontinuidades temporales & 0 \\
        Estados ADS inválidos & 0 \\
        RMS & \SI{0.115}{\micro\volt} \\
        Diagnóstico & Prueba válida \\
        \hline
    \end{tabular}
\end{table}

\begin{figure}[!htbp]
    \centering
    \validacioninclude{width=0.82\textwidth}{images/validacion/fig_01_shorted_inputs_timeseries.png}
    \caption{Señal temporal registrada en la prueba diagnóstica con entradas cortocircuitadas. La amplitud reducida y la ausencia de discontinuidades permiten verificar que la ruta ADS1299--firmware--Python funciona correctamente antes de conectar el montaje experimental.}
    \label{fig:validacion_shorted_inputs}
\end{figure}

Una vez comprobada la ruta de adquisición, se compararon diferentes configuraciones físicas de electrodos y modos de operación del ADS1299. Para esta validación se utilizaron ocho capturas específicas, agrupadas en dos montajes: \texttt{ear\_eeg\_ch1\_only} y \texttt{fp1\_fp2\_ch1\_only}. Para evitar que la tabla quede sobrecargada con nombres de carpeta demasiado largos, cada captura se identifica mediante un código corto. La correspondencia con las carpetas originales se mantiene por la marca temporal y la condición experimental.

\begin{table}[!htbp]
    \centering
    \small
    \caption{Capturas utilizadas para validar el montaje experimental.}
    \label{tab:validacion_capturas_montaje}
    \begin{tabular}{c p{0.27\textwidth} p{0.25\textwidth} p{0.28\textwidth}}
        \hline
        \textbf{ID} & \textbf{Montaje} & \textbf{Condición} & \textbf{Carpeta de captura} \\
        \hline
        M1 & Ear-EEG CH1 & Reposo quieto & \texttt{20260523-195752} \\
        M2 & Ear-EEG CH1 & Ojos abiertos & \texttt{20260523-200925} \\
        M3 & Ear-EEG CH1 & Ojos cerrados & \texttt{20260523-201055} \\
        M4 & Ear-EEG CH1 & Movimiento de mandíbula & \texttt{20260523-201321} \\
        M5 & Fp1--Fp2 CH1 & Reposo quieto & \texttt{20260523-202120} \\
        M6 & Fp1--Fp2 CH1 & Ojos abiertos & \texttt{20260523-202208} \\
        M7 & Fp1--Fp2 CH1 & Ojos cerrados & \texttt{20260523-202323} \\
        M8 & Fp1--Fp2 CH1 & Parpadeo/frente & \texttt{20260523-202451} \\
        \hline
    \end{tabular}
\end{table}

Las carpetas completas asociadas a estos identificadores son, respectivamente: \texttt{ear\_eeg\_ch1\_only\_still\_30s}, \texttt{ear\_eeg\_ch1\_only\_eyes\_open\_60s}, \texttt{ear\_eeg\_ch1\_only\_eyes\_closed\_60s}, \texttt{ear\_eeg\_ch1\_only\_jaw\_movement\_30s}, \texttt{fp1\_fp2\_ch1\_only\_quiet\_30s}, \texttt{fp1\_fp2\_ch1\_only\_eyes\_open\_60s}, \texttt{fp1\_fp2\_ch1\_only\_eyes\_closed\_60s} y \texttt{fp1\_fp2\_ch1\_only\_forehead\_blink\_artifact\_30s}. De este modo, la tabla conserva la trazabilidad de las capturas sin comprometer la legibilidad de la página.

La primera lectura de estas capturas se realizó en el dominio temporal. El objetivo fue comprobar que las señales quedaban registradas de forma continua en todas las condiciones y que las diferencias observadas procedían del montaje o del estado experimental, no de fallos de transporte. Las Figuras~\ref{fig:validacion_montaje_temporales_ear} y~\ref{fig:validacion_montaje_temporales_fp1fp2} resumen las señales temporales de las ocho capturas usadas para comparar montajes, separando el montaje auricular y el montaje frontal para mejorar la legibilidad.

\begin{figure}[!htbp]
    \centering
    \validacioninclude{width=0.96\textwidth}{images/validacion/fig_02_mounting_timeseries_grid_ear.png}
    \caption{Comparación temporal de las capturas realizadas con el montaje \texttt{ear\_eeg\_ch1\_only}. La figura agrupa reposo, ojos abiertos, ojos cerrados y movimiento de mandíbula, permitiendo comprobar la continuidad de adquisición y la respuesta del montaje ante distintas condiciones.}
    \label{fig:validacion_montaje_temporales_ear}
\end{figure}

\begin{figure}[!htbp]
    \centering
    \validacioninclude{width=0.96\textwidth}{images/validacion/fig_02_mounting_timeseries_grid_fp1fp2.png}
    \caption{Comparación temporal de las capturas realizadas con el montaje \texttt{fp1\_fp2\_ch1\_only}. La figura agrupa reposo, ojos abiertos, ojos cerrados y parpadeo/frente, mostrando una mayor sensibilidad a artefactos frontales y de contacto.}
    \label{fig:validacion_montaje_temporales_fp1fp2}
\end{figure}

El montaje frontal Fp1--Fp2 permitió observar actividad y artefactos claros, especialmente asociados a parpadeos y movimiento frontal, pero mostró mayor sensibilidad al contacto, al movimiento y al ruido común. Por el contrario, el montaje ear-EEG centrado en CH1 proporcionó una base más estable para obtener ventanas útiles, reduciendo la influencia de canales no usados y facilitando la aplicación posterior del \textit{quality gate}.

La configuración final quedó fijada como \texttt{ADS\_DIAGNOSTIC\_MODE=5}, modo \texttt{bias\_ch1\_only\_loff\_off} y montaje \texttt{ear\_eeg\_ch1\_only}. En este modo, CH1 permanece como canal principal de análisis, el BIAS se deriva de CH1P y CH1N, la detección \textit{lead-off} se mantiene desactivada y los canales CH2--CH4 se conservan por compatibilidad con el contrato de datos, pero no se emplean como canales activos para la cadena de sonificación.

Para justificar esta elección no se utilizó una única métrica, sino un conjunto de indicadores complementarios. El RMS permite comparar la amplitud típica de la señal; el pico a pico permite detectar transitorios grandes; la componente relativa de 50 Hz indica sensibilidad al ruido de red; la fracción de artefactos resume la proporción de ventanas contaminadas; y el \textit{quality score} integra varios de estos criterios en una puntuación de calidad útil para decidir si una ventana puede alimentar la sonificación. Las Figuras~\ref{fig:validacion_montaje_metricas} y~\ref{fig:validacion_montaje_quality_score} agrupan estas evidencias.

\begin{figure}[!htbp]
    \centering
    \begin{minipage}{0.48\textwidth}
        \centering
        \validacioninclude{width=\textwidth}{images/validacion/fig_02_mounting_rms_comparison.png}
        \small (a) RMS mediano
    \end{minipage}
    \hfill
    \begin{minipage}{0.48\textwidth}
        \centering
        \validacioninclude{width=\textwidth}{images/validacion/fig_02_mounting_ptp_comparison.png}
        \small (b) Amplitud pico a pico
    \end{minipage}

    \vspace{0.25cm}

    \begin{minipage}{0.48\textwidth}
        \centering
        \validacioninclude{width=\textwidth}{images/validacion/fig_02_mounting_50hz_comparison.png}
        \small (c) Componente relativa de 50 Hz
    \end{minipage}
    \hfill
    \begin{minipage}{0.48\textwidth}
        \centering
        \validacioninclude{width=\textwidth}{images/validacion/fig_02_mounting_artifact_fraction.png}
        \small (d) Fracción de artefactos
    \end{minipage}
    \caption{Comparación de montajes mediante métricas de amplitud, ruido de red y artefactos. El análisis conjunto permite valorar la estabilidad relativa de cada configuración y detectar condiciones con mayor sensibilidad a transitorios o contaminación.}
    \label{fig:validacion_montaje_metricas}
\end{figure}

\begin{figure}[!htbp]
    \centering
    \validacioninclude{width=0.82\textwidth}{images/validacion/fig_02_mounting_quality_score.png}
    \caption{Comparación del \textit{quality score} entre montajes y condiciones. La puntuación resume la validez relativa de las ventanas para alimentar el análisis espectral y la sonificación.}
    \label{fig:validacion_montaje_quality_score}
\end{figure}

La configuración \texttt{ear\_eeg\_ch1\_only} con \texttt{bias\_ch1\_only\_loff\_off} se seleccionó como montaje final porque ofreció una base más estable para obtener ventanas útiles y aplicar posteriormente el \textit{quality gate}. Los resultados no implican que el montaje final elimine todos los artefactos biológicos o ambientales. Su utilidad reside en que proporciona una señal suficientemente estable para trabajar con ventanas útiles y permite que el sistema de calidad descarte o atenúe las ventanas contaminadas antes de la sonificación, sin actuar de manera constante ni brusca.

Por tanto, la elección del montaje final se justifica por la estabilidad temporal, la ausencia de discontinuidades, el estado ADS válido, la menor influencia de canales no usados y su compatibilidad con el \textit{quality gate} utilizado en la validación final. La sesión final de laboratorio, analizada posteriormente en la subsección de sonificación y resultados, se realizó con esta configuración. En consecuencia, las limitaciones observadas en la señal real no se atribuyen a pérdidas de transporte ni a fallos de estado del ADS1299, sino principalmente al montaje bioeléctrico, al contacto de electrodos, al ruido de red y a artefactos fisiológicos o externos.
